%==============================================================
% BAB IV SKALABILITAS, KELAYAKAN, DAN RENCANA PENGEMBANGAN
%==============================================================

\section{Skalabilitas, Kelayakan, dan Rencana Pengembangan}

\subsection{Arsitektur Skalabilitas Tiga Dimensi (Teknis, Konten, dan Jangkauan)}

Agar dapat tumbuh dari purwarupa fungsional menjadi layanan yang melayani ribuan pengguna di seluruh Indonesia tanpa membebani infrastruktur, TropiCare dirancang dengan skalabilitas pada tiga dimensi utama:

\begin{enumerate}
    \item \textbf{Skalabilitas Teknis (\textit{Stateless \& Modular}):} 
    Endpoint skoring pada backend FastAPI bersifat \textit{stateless} sehingga dapat direplikasi secara horizontal dengan mudah. Arsitektur menggunakan pola \textit{Ports and Adapters} (heksagonal): mesin skoring aturan deterministik maupun penyedia model AI (Gemma 4 / MedGemma via OpenRouter, vLLM mandiri, atau Gemini) dipasang pada lapisan komposisi (\textit{composition root}) dan dapat diganti hanya melalui konfigurasi tanpa mengubah logika bisnis. Selain itu, deteksi 11 tanda bahaya darurat dijalankan sepenuhnya di peramban pasien (biaya komputasi server nol), dan pratinjau skor selama pengisian kuis memakai mesin aturan lokal serta \textit{cache} memori 24 jam sehingga pemanggilan LLM hanya terjadi satu kali di akhir sesi skrining non-darurat.

    \item \textbf{Skalabilitas Konten Klinis (\textit{Data-Driven JSON}):} 
    Seluruh pertanyaan, opsi jawaban, aturan peringatan obat, dan bobot penyakit disimpan dalam berkas konfigurasi JSON terstruktur yang terpisah dari kode antarmuka React. Dokter dan statistisi dapat menambahkan penyakit baru, memperbarui bobot gejala, atau menyesuaikan profil endemisitas wilayah secara langsung tanpa perlu memodifikasi kode tampilan frontend. Setiap perubahan teks medis dan aturan klinis dikendalikan melalui persetujuan kolaborator dokter.

    \item \textbf{Skalabilitas Jangkauan \& Surveilans Berprivasi:} 
    Ekspansi layanan dirancang melalui lima tahap berurutan: \textbf{(1) Validasi Klinis} $\rightarrow$ \textbf{(2) Pilot Klinik \& Puskesmas} $\rightarrow$ \textbf{(3) TropiCare Plus (Pemantauan PTM)} $\rightarrow$ \textbf{(4) White-label untuk LSM Kesehatan} $\rightarrow$ \textbf{(5) API Mitra \& Surveilans Dinas Kesehatan}. Untuk mendukung Dinas Kesehatan tanpa melanggar privasi warga, statistik wilayah disajikan murni dari agregat tanpa identitas langsung dengan aturan privasi ketat: \textbf{kelompok agregat dengan jumlah kasus $<10$ disembunyikan} guna mencegah pelacakan balik individu.
\end{enumerate}

\subsection{Analisis Kelayakan Finansial (\textit{Feasibility}) dan Ekonomi Unit}

Keunggulan utama desain arsitektur ringan TropiCare adalah biaya operasional bulanan yang sangat rendah, sehingga layanan ini dapat bertahan hidup secara mandiri tanpa harus menunggu anggaran pengadaan pemerintah yang panjang. Rincian perhitungan biaya infrastruktur bulanan dan model keberlanjutan disajikan pada Tabel~\ref{tab:feasibility}.

\begin{table}[H]
\centering
\footnotesize
\caption{Rincian Biaya Operasional Bulanan dan Skema Monetisasi Berkelanjutan TropiCare}\label{tab:feasibility}
\vspace{3pt}
\begin{tabularx}{\textwidth}{p{3.4cm} p{2.8cm} X}
\toprule
\textbf{Komponen Biaya / Layanan} & \textbf{Nilai Finansial} & \textbf{Dasar Perhitungan Teknis \& Model Bisnis} \\
\midrule
\multicolumn{3}{l}{\textbf{A. Struktur Biaya Operasional Aplikasi (\textit{Monthly Burn Rate})}} \\
\midrule
\textbf{Server \& Basis Data\newline(\textit{Railway Hobby Plan})} & \textbf{Rp90.000 / bulan} & Menjalankan kontainer backend FastAPI \textit{stateless} dan basis data PostgreSQL (menyimpan 13.930 mirror faskes, riwayat akun opsional, dan agregat statistik). \\
\midrule
\textbf{Inferensi AI\newline(\textit{OpenRouter API})} & \textbf{Rp180.280 / bulan} & Melayani panggilan tunggal terstruktur (\textit{single-call JSON}) model keluarga Gemma + pengenalan obat. Dengan $\sim 1.550$ token/skrining dan \textit{cache} 24 jam, anggaran ini mencukupi $\sim \mathbf{15.000\text{--}20.000}$ skrining penuh/bulan ($\approx \text{Rp}9\text{--}\text{Rp}12$ per pasien). \\
\midrule
\textbf{Domain \& DNS Resmi} & \textbf{Rp25.000 / bulan} & Biaya tahunan nama domain resmi (\text{Rp}300.000/tahun) dibagi rata per bulan. \\
\midrule
\textbf{Total Biaya Server} & \textbf{Rp295.280 / bulan} & \textbf{Titik Impas (\textit{Break-Even}): Cukup 1 langganan klinik mitra (Rp400.000/bulan) sudah menutup seluruh biaya server sebulan penuh.} \\
\midrule
\multicolumn{3}{l}{\textbf{B. Skema Keberlanjutan Pendapatan (Skrining Inti Masyarakat Tetap Gratis)}} \\
\midrule
\textbf{1. Dasbor Faskes\newline(Klinik \& Puskesmas)} & \textbf{Rp250.000--Rp400.000\newline/ bulan / faskes} & Layanan B2B pra-triase loket lewat pemindaian QR dan rekapitulasi keluhan pasien datang untuk memangkas antrean administrasi. \\
\midrule
\textbf{2. Lisensi API \& Kemitraan Surveilans} & \textbf{Rp15.000.000--Rp40.000.000 / tahun} & Akses agregat tren penyakit per wilayah tanpa identitas langsung (kelompok $<10$ disembunyikan) bagi LSM/NGO eliminasi penyakit tropis, institusi riset, atau kemitraan CSR kesehatan. \\
\midrule
\textbf{3. TropiCare Plus\newline(Pantau PTM Mandiri)} & \textbf{Rp15.000 / bulan\newline(Opsional)} & Fitur tambahan bagi penderita hipertensi atau diabetes untuk memantau grafik tren tekanan darah dan gula darah jangka panjang. \\
\bottomrule
\end{tabularx}
\end{table}

\subsection{Peta Jalan Pengembangan, Mitigasi Risiko, dan Indikator Kinerja (KPI)}

Pelaksanaan pengembangan selama 12 bulan dari tahap validasi hingga integrasi ekosistem, beserta matriks mitigasi risiko dan indikator kinerja utama (KPI), dirangkum pada Tabel~\ref{tab:roadmap} dan Tabel~\ref{tab:kpi}.

\begin{table}[H]
\centering
\footnotesize
\caption{Peta Jalan Pengembangan 12 Bulan dan Matriks Mitigasi Risiko Utama}\label{tab:roadmap}
\vspace{3pt}
\begin{tabularx}{\textwidth}{p{2.5cm} p{3.1cm} X p{3.8cm}}
\toprule
\textbf{Tahap \& Periode} & \textbf{Fokus \& Risiko Utama} & \textbf{Aktivitas Eksekusi \& Mitigasi Konkret} & \textbf{Luaran Kunci (\textit{Milestone})} \\
\midrule
\textbf{Tahap 1\newline(Bulan 1--3)} & Validasi Klinis \& Mitigasi Risiko Halusinasi AI & Uji buta 100 kartu kasus dokter; kalibrasi bobot aturan JSON bersama statistisi; pemutusan 11 tanda bahaya langsung di browser tanpa AI (\textit{safety gate}); uji kemudahan pakai (SUS). & Akurasi triase $>92\%$ tanpa kasus darurat terlewat; skor SUS $>80$ (\textit{Grade A}); 304 kasus uji backend lulus. \\
\midrule
\textbf{Tahap 2\newline(Bulan 4--6)} & Pilot Klinik/FKTP \& Mitigasi Beban Loket & Pilot di 10 klinik pratama dan Puskesmas wilayah endemis; penerapan pemindaian QR rangkuman tanpa ketik ulang (\textit{zero-retyping}); kurasi nomor WhatsApp aktif faskes. & Operasional di 10 faskes pilot; biaya server tertutup penuh dari langganan mitra awal (\textit{self-sustaining}). \\
\midrule
\textbf{Tahap 3\newline(Bulan 7--9)} & TropiCare Plus, Regulasi \& Privasi Data & Peluncuran modul grafik Pantau PTM; pendaftaran \textit{Regulatory Sandbox} Kemenkes RI \citep{UUKesehatan2023,KemenkesDTO2024}; audit kepatuhan UU PDP No. 27/2022 (Zero-NIK, lokasi tetap di perangkat) \citep{UUPDP2022}. & Terdaftar di \textit{Regulatory Sandbox} Kemenkes; modul Pantau PTM aktif; jembatan rujukan Mobile JKN teruji. \\
\midrule
\textbf{Tahap 4--5\newline(Bulan 10--12)} & White-Label LSM \& API Mitra Surveilans & Kemitraan white-label dengan LSM pemberantasan penyakit tropis (TBC, Kusta, NTD); pembukaan API agregat wilayah bagi Dinkes (kelompok $<10$ disembunyikan); persiapan pemetaan standar FHIR R4. & Adopsi di 100+ faskes; dasbor surveilans agregat wilayah beroperasi; kesiapan interoperabilitas nasional. \\
\bottomrule
\end{tabularx}
\end{table}

\begin{table}[H]
\centering
\footnotesize
\caption{Indikator Kinerja Utama (KPI) Keberhasilan Implementasi TropiCare}\label{tab:kpi}
\vspace{3pt}
\begin{tabularx}{\textwidth}{p{1.4cm} p{4.2cm} p{2.2cm} X}
\toprule
\textbf{Kode} & \textbf{Indikator Kinerja} & \textbf{Target} & \textbf{Metode Pengukuran} \\
\midrule
\textbf{KPI-1} & Akurasi Triase terhadap Dokter & $\mathbf{>92\%}$ & Uji kesesuaian rute triase pada 100 kartu kasus klinis dokter, dengan syarat tidak ada satu pun kondisi darurat yang terlewat. \\
\midrule
\textbf{KPI-2} & Efisiensi Waktu Registrasi Loket & $\mathbf{\ge 40\%}$ & Pemangkasan waktu wawancara awal di loket faskes melalui pemindaian QR rangkuman dibandingkan pendaftaran manual. \\
\midrule
\textbf{KPI-3} & Kepatuhan Puasa Pra-Laboratorium & $\mathbf{>85\%}$ & Proporsi pasien skrining jalur PTM yang datang ke FKTP memenuhi syarat puasa 10--12 jam berkat instruksi pada hasil skrining. \\
\midrule
\textbf{KPI-4} & Kecepatan Deteksi Tanda Bahaya & $\mathbf{<500\text{ ms}}$ & Waktu respon peramban dari ketukan tanda bahaya hingga layar Rute Merah dan tombol 119 tampil di ponsel pasien. \\
\midrule
\textbf{KPI-5} & Skor Kemudahan Pakai (\textit{SUS}) & $\mathbf{>80}$ (\textit{Grade A}) & Evaluasi kuesioner standar \textit{System Usability Scale} pada masyarakat umum dan petugas faskes. \\
\bottomrule
\end{tabularx}
\end{table}
